\section{Numerical results}
\label{sec:results}

We first compare the exact noisy fidelity of the validation-selected
circuits across all targets. We then include the classical and circuit
resources required to obtain those endpoints. The native-gate pilot tests
how much of the resulting interpretation survives a different placement
of the noise windows.

\subsection{Fidelity across target families}

\Cref{fig:fidelity} shows the complete logical-noise comparison. The
top-down workflow attains the highest selected-endpoint fidelity on five
of nine targets. Its mean fidelity is 92.260\%, compared with 91.019\% for
RobustState. The paired mean difference is 1.241 percentage points, with
five wins and four losses. These are descriptive comparisons of the
selected circuits, without an estimate of optimizer-seed variability.

\begin{figure*}[t]
 \centering
 \includegraphics[width=\textwidth]{figures/fidelity.pdf}
 \caption{\textbf{Selected fidelity under logical-gate noise.}
 (a) Exact density-matrix fidelity for the nine fixed six-qubit targets,
 optimizer seed 300, and the synthetic Pauli profile in
 \cref{eq:jump_law}. Each value refers to a circuit selected by independent
 validation, with archived top-down selections retained. Fontana excludes
 the three complex Haar targets. Bold entries mark the highest fidelity
 for each target. (b) Paired top-down minus RobustState differences in
 percentage points; blue and orange denote positive and negative values.
 The top-down workflow has the highest fidelity on five targets.
 Exact endpoint scores have no trajectory sampling error, while
 optimizer-seed variability remains unmeasured.}
 \label{fig:fidelity}
\end{figure*}

The six real-valued targets give the strongest results for the top-down
method. Their mean fidelity is 95.885\%, compared with 94.153\% for
RobustState. The top-down method wins on Gaussian, ferromagnetic and
critical Ising, and bond-two MPS. RobustState is better on paramagnetic
Ising and bond-three MPS by 0.237 and 0.249 percentage points, respectively.
This division describes the tested fixture. It does not isolate real
amplitudes, entanglement, or another target property as the cause.

The complex Haar targets expose a different regime. RobustState is more
accurate on Haar 1 and Haar 3, with gaps of 1.132 and 0.551 percentage
points. On Haar 2, the top-down method instead gains 2.463 percentage
points. Thus even these three targets do not produce a uniform ranking.
Their top-down fidelities remain between 84.811\% and 85.250\%, below the
95.272--96.472\% range of the six real targets.

Fontana is supported on the six real targets only. On that common subset,
the top-down workflow wins in all six cases, with a mean advantage of
3.875 percentage points. Fontana's six-target mean must be interpreted on
this subset, rather than compared with nine-target means. NACL has lower
fidelity on all nine targets at the declared 50-step structural-search
budget. The result characterizes that finite search and does not establish
the performance of a converged NACL implementation.

The top-down workflow also exceeds both internal controls in fidelity on
every target. The mean improvements are 4.611 percentage points over
fixed VQSP and 2.929 over the unpruned layerwise control. The controls
differ in candidate grids and optimization policies, so these results
support the complete preparation workflow without assigning the gain
to pruning or noisy refinement alone.

\subsection{Classical cost and circuit resources}

\Cref{fig:pareto} combines selected fidelity with full construction and
training CPU. The top-down workflow is Pareto-efficient on seven of nine
targets, both against the external adaptations alone and when the internal
controls are included. Across the complete fixture it uses 285.17 CPU
seconds, compared with 256.52 for RobustState, an increase of 11.17\%.
The modest aggregate fidelity advantage therefore has an aggregate
classical cost.

\begin{figure*}[t]
 \centering
 \includegraphics[width=\textwidth]{figures/pareto.pdf}
 \caption{\textbf{Fidelity versus complete construction and training cost.}
 Each panel compares the selected endpoints for one target under the
 logical-gate noise model. CPU includes every visited candidate and the
 historical work of reused searches; it excludes validation and final
 assessment. The horizontal axes are logarithmic. The marked Pareto
 frontiers use all supported methods, including controls. Connections
 identify the discrete frontier and do not represent an interpolated
 performance curve. The top-down workflow is nondominated on seven
 targets; RobustState dominates it on Haar 1 and Haar 3.}
 \label{fig:pareto}
\end{figure*}

The cost separates sharply by target subset. On the six real targets,
top-down construction and training take 103.81 CPU seconds, compared with
170.92 for RobustState, a reduction of 39.27\%. Top-down is faster on each
of these targets. The two small fidelity losses to RobustState therefore
remain fidelity--runtime trade-offs. Fontana is substantially cheaper
still, taking 31.37 CPU seconds on the same six targets. The top-down
method's paired fidelity advantage over Fontana costs 3.31 times as much
classical computation.

RobustState dominates top-down on Haar 1 and Haar 3 in both fidelity and
CPU. It also uses fewer logical two-qubit gates and less two-qubit depth.
On Haar 2, the top-down fidelity gain costs 53.16 rather than 29.35 CPU
seconds and requires 38 rather than 20 logical two-qubit gates. Higher
fidelity on this target is therefore obtained with more classical work
and a larger logical circuit.

The fixed VQSP control is cheaper overall, at 115.88 CPU seconds, while
the layerwise control takes 354.24 seconds. The latter has both lower
mean fidelity and greater aggregate CPU than top-down. Aggregate values
do not replace the per-target Pareto comparisons. Full fidelity and
resource tables are included in the appendices.

These costs charge the complete method histories. The extension campaign
itself took 374.21 seconds of wall time and performed 966.96 CPU seconds
of fresh construction and training. Including historical work, the
charged total across all methods is 1,753.68 CPU seconds. The campaign
wall time is consequently not a measure of any single method's cost.

\subsection{Native-gate noise}

The logical comparison assigns one arity-dependent window to every
original elementary gate. To examine the effect of native implementation,
we also report the matched three-target pilot in \cref{tab:native}.
Its circuits are trained, selected, and assessed with noise after a frozen
$R_Z$/SX/$X$/CX schedule. The top-down method reuses the same noiseless
construction and pruning B checkpoints, followed by native refinement.
All historical construction work is charged.

\input{tables/native_results}

The native comparison is less favorable to top-down. On Gaussian,
Fontana reaches 94.144\% fidelity at 19.66 CPU seconds, compared with
92.891\% at 26.17 seconds for top-down. On Haar 1, RobustState reaches
82.090\% at 52.24 seconds, compared with 70.220\% at 79.84 seconds.
Both competitors also use fewer CX gates on the corresponding target.
These are direct counterexamples to a uniform native advantage.

Critical Ising retains a fidelity--runtime trade-off. Top-down reaches
90.062\%, compared with 90.334\% for RobustState, while taking 33.46
rather than 52.29 CPU seconds. Its 36.0\% CPU saving accompanies a
0.272-percentage-point fidelity loss and 26 rather than 10 CX gates.
Fontana is faster than either at 9.65 seconds, with 88.005\% fidelity.

After native selection was locked, a separate diagnostic evaluated the
earlier logically trained parameters on the same frozen native schedules
at the selected depths. Relative to those parameters, native refinement
changed fidelity by $-0.021$, $-0.026$, and $+0.020$ percentage points for
Gaussian, critical Ising, and Haar 1. The small changes do not show a
substantial benefit from changing the training noise location in this
pilot. They also should not be confused with the larger differences from
earlier numerically compiled deployment circuits, which involve schedule
changes and, for Haar 1, a different selected depth.

Native refinement nevertheless improves upon the pruned starting circuits
by 0.974, 4.446, and 3.842 percentage points, respectively. The absence of
a further gain over logical training on the matched schedules is a
different observation from the gain obtained by refining the pruned input.
The pilot uses one seed and three targets, with 32,768-path top-down
validation and the competitors' retained 8,192-path selections. It limits
the interpretation of the logical results without establishing a general
native-gate ranking.
